\documentclass[a4paper,12pt]{jsarticle}
\usepackage[margin=25mm]{geometry}
\usepackage{amsmath,amssymb}
\usepackage{enumitem}

\begin{document}

\begin{center}
{\LARGE 数I・A 共通テスト本番レベル（図形のみ）解答}
\end{center}

\vspace{1em}

\section*{第1問}
\begin{enumerate}[label=\arabic*.]
  \item 余弦定理より
  \[
  BC^2=7^2+9^2-2\cdot7\cdot9\cos120^\circ
  =49+81+63=193.
  \]
  よって
  \[
  BC=\sqrt{193}.
  \]
  \item 面積
  \[
  S=\frac12\cdot7\cdot9\sin120^\circ
  =\frac{63}{2}\cdot\frac{\sqrt3}{2}
  =\frac{63\sqrt3}{4}.
  \]
  \item 中線の長さ公式
  \[
  AM=\frac12\sqrt{2AB^2+2AC^2-BC^2}
  =\frac12\sqrt{2\cdot49+2\cdot81-193}
  =\frac{\sqrt{67}}{2}.
  \]
  \item 正弦定理 \( \dfrac{BC}{\sin A}=2R \) より
  \[
  R=\frac{BC}{2\sin120^\circ}
  =\frac{\sqrt{193}}{\sqrt3}
  =\frac{\sqrt{579}}{3}.
  \]
\end{enumerate}

\section*{第2問}
\begin{enumerate}[label=\arabic*.]
  \item 弦の公式 \(AB=2R\sin\frac{\theta}{2}\)（\(\theta=\angle AOB\)）より
  \[
  8\sqrt2=16\sin\frac{\theta}{2}
  \Rightarrow \sin\frac{\theta}{2}=\frac{\sqrt2}{2}.
  \]
  小さい方より
  \[
  \theta=90^\circ.
  \]
  \item 弧長
  \[
  \frac{90}{360}\cdot2\pi\cdot8=4\pi.
  \]
  \item 扇形面積
  \[
  \frac{90}{360}\cdot\pi\cdot8^2=16\pi.
  \]
  \item 弓形面積
  \[
  =\text{扇形}-\triangle AOB.
  \]
  \[
  [\triangle AOB]=\frac12\cdot8\cdot8\sin90^\circ=32.
  \]
  よって
  \[
  16\pi-32.
  \]
\end{enumerate}

\section*{第3問}
\begin{enumerate}[label=\arabic*.]
  \item 面積
  \[
  S=AB\cdot AD\sin60^\circ
  =10\cdot6\cdot\frac{\sqrt3}{2}
  =30\sqrt3.
  \]
  \item 三角形 \(ABD\) に余弦定理を用いて
  \[
  BD^2=10^2+6^2-2\cdot10\cdot6\cos60^\circ
  =100+36-60=76.
  \]
  よって
  \[
  BD=2\sqrt{19}.
  \]
  \item 対角線は平行四辺形を等面積の4三角形に分けるので
  \[
  [\triangle AOD]=\frac14\cdot30\sqrt3=\frac{15\sqrt3}{2}.
  \]
  \item
  \[
  [\triangle ABD]=\frac12\cdot30\sqrt3=15\sqrt3.
  \]
  点 \(A\) から直線 \(BD\) への距離を \(d\) とすると
  \[
  15\sqrt3=\frac12\cdot(2\sqrt{19})\cdot d
  \Rightarrow d=\frac{15\sqrt3}{\sqrt{19}}
  =\frac{15\sqrt{57}}{19}.
  \]
\end{enumerate}

\section*{第4問}
座標を
\[
A(0,0,0),\ B(6,0,0),\ D(0,6,0),\ E(0,0,9),\ G(6,6,9)
\]
とおく。
\begin{enumerate}[label=\arabic*.]
  \item
  \[
  AG=\sqrt{6^2+6^2+9^2}
  =\sqrt{153}=3\sqrt{17}.
  \]
  \item \(AG\) の底面への正射影は \(AC\) で
  \[
  AC=\sqrt{6^2+6^2}=6\sqrt2.
  \]
  よって
  \[
  \tan\theta=\frac{9}{6\sqrt2}=\frac{3\sqrt2}{4}.
  \]
  \item
  \[
  [\triangle AEG]
  =\frac12\cdot AE\cdot(\text{点 }G\text{ から直線 }AE\text{ への距離})
  =\frac12\cdot9\cdot6\sqrt2
  =27\sqrt2.
  \]
  \item
  \[
  \overrightarrow{BD}=(-6,6,0),\quad \overrightarrow{BG}=(0,6,9).
  \]
  法線ベクトルを
  \[
  \overrightarrow{n}=\overrightarrow{BD}\times\overrightarrow{BG}=(3,3,-2)
  \]
  （比例ベクトル）とすると、平面 \(BDG\) の方程式は
  \[
  3x+3y-2z-18=0.
  \]
  よって点 \(A\) からこの平面までの距離は
  \[
  \frac{|3\cdot0+3\cdot0-2\cdot0-18|}{\sqrt{3^2+3^2+(-2)^2}}
  =\frac{18}{\sqrt{22}}
  =\frac{9\sqrt{22}}{11}.
  \]
\end{enumerate}

\end{document}
